\setcounter{aufgabe}{31}
\hypertarget{e4}{}\einheitenkopf{Einheit 4 von 5 · Gleichung bestimmen}
\zweigzeile{Hier lernst du, die Gleichung einer Geraden aus dem Graphen, aus Steigung und Punkt oder aus zwei Punkten zu bestimmen und den Schnittpunkt zweier Geraden zu berechnen · neu in diesem Jahr · P10 · baut auf: $m$ und $n$ ablesen (Einheit 2), Punktprobe und Nullstelle (Einheit 3), lineare Gleichungen, Terme zusammenfassen}

\begin{aufgabe}{\textbf{Ich kann die Gleichung einer Geraden am Graphen ablesen.} Lies $m$ und $n$ ab und gib die Gleichung an.}

\begin{ksys}[xmin=-4,xmax=4,ymin=-5,ymax=6,ablesen]
\gerade[3.5]{1}{-3}{g}\gerade[-1]{-3}{2}{h}\gerade[-3.5]{2}{4}{k}
\end{ksys}

\begin{geruest}
\gz{a}{Gerade $g$}{\feld{m}\feld{n}\feldl{y}}
\gz{b}{Gerade $h$}{\feld{m}\feld{n}\feldl{y}}
\gz{c}{Gerade $k$}{\feld{m}\feld{n}\feldl{y}}
\end{geruest}
\end{aufgabe}

\begin{aufgabe}{\textbf{Ich kann die Gleichung aus der Steigung und einem Punkt bestimmen.} Setze $m$ und die Koordinaten des Punktes in $y = m \cdot x + n$ ein, berechne $n$ und gib die Gleichung an.}
\beispiel{m = 2,\ P(3\,|\,5):\quad 5 &= 2 \cdot 3 + n \\ 5 &= 6 + n &&\mid -6 \\ n &= -1 && y = 2x - 1}
\begin{gleichungsraster}[3]
\gl{m = 2,\ P(1\,|\,5)} & \gl{m = 3,\ P(2\,|\,4)} \\
\gl{m = -1,\ P(3\,|\,1)} & \gl{m = -2,\ P(-3\,|\,4)} \\
\gl{m = 0{,}5,\ P(-4\,|\,3)} & \\
\end{gleichungsraster}
\end{aufgabe}

\begin{aufgabe}{\textbf{Ich kann die Gleichung einer Geraden durch zwei Punkte bestimmen.} Berechne zuerst $m$: Differenz der $y$-Werte geteilt durch die Differenz der $x$-Werte, beide Male in derselben Reihenfolge. Berechne dann $n$ und gib die Gleichung an.}
\beispiel{A(2\,|\,3),\ B(4\,|\,9):\quad m &= (9 - 3) : (4 - 2) = 6 : 2 = 3 \\ 3 &= 3 \cdot 2 + n &&\mid -6 \\ n &= -3 && y = 3x - 3}
\begin{gleichungsraster}[3]
\gl{A(0\,|\,2),\ B(2\,|\,8)} & \gl{A(0\,|\,{-1}),\ B(4\,|\,7)} \\
\gl{A(1\,|\,6),\ B(4\,|\,15)} & \gl{A(-1\,|\,7),\ B(2\,|\,{-2})} \\
\end{gleichungsraster}
\end{aufgabe}

\begin{aufgabe}{\textbf{Ich kann eine Gerade durch zwei Punkte zeichnen.} Trage die Punkte ein und zeichne die Gerade durch beide Punkte. Zeichne a) und b) links, c) rechts.}
\begin{teile}
\teil $A(0\,|\,1)$ und $B(2\,|\,5)$
\teil $C(-3\,|\,4)$ und $D(3\,|\,{-2})$
\teil Zeichne die Gerade durch $E(-2\,|\,4)$ und $F(4\,|\,1)$. Bestimme $m$ und $n$ und gib die Gleichung an. (P10 2025) \\ \feld{m} \feld{n} \feldl{y}
\end{teile}

\begin{ksys}[xmin=-4,xmax=4,ymin=-3,ymax=6]
\end{ksys}\ksysabstand\begin{ksys}[xmin=-3,xmax=5,ymin=-3,ymax=6]
\end{ksys}

\end{aufgabe}

\begin{aufgabe}{\textbf{Ich kann den Schnittpunkt zweier Geraden berechnen.} Setze die beiden Terme gleich, löse nach $x$ auf und berechne $y$ mit einer der beiden Gleichungen.}
\beispiel{f(x) = 3x - 2,\ g(x) = x + 4:\quad 3x - 2 &= x + 4 &&\mid -x \\ 2x - 2 &= 4 &&\mid +2 \\ 2x &= 6 &&\mid :2 \\ x &= 3 && y = 3 + 4 = 7,\ S(3\,|\,7)}
\begin{gleichungsraster}[4]
\gl{f(x) = 3x + 2,\ g(x) = x + 6} & \gl{f(x) = 4x - 3,\ g(x) = 2x + 5} \\
\gl{f(x) = 5x - 4,\ g(x) = 2x + 2} & \gl{f(x) = 2x + 7,\ g(x) = 5x - 2} \\
\end{gleichungsraster}
\end{aufgabe}

\begin{aufgabe}{\textbf{Ich kann den Schnittpunkt zweier Geraden berechnen – weiter.} Berechne den Schnittpunkt.}
\setcounter{teil}{4}
\begin{gleichungsraster}[4]
\gl{f(x) = -x + 7,\ g(x) = 2x - 2} & \gl{f(x) = 0{,}5x + 1,\ g(x) = -x + 4} \\
\gl{f(x) = -2x + 3,\ g(x) = 2x - 3} & \\
\end{gleichungsraster}
\end{aufgabe}

\begin{aufgabe}{\textbf{Ich finde den Fehler bei der Steigung aus zwei Punkten.} Mia bestimmt die Gerade durch $P(2\,|\,5)$ und $Q(4\,|\,11)$ so:}
\rechnung{m &= (5 - 11) : (4 - 2) = -3 \\ 5 &= -3 \cdot 2 + n \\ n &= 11 && y = -3x + 11}
\begin{teile}
\teil Finde den Fehler, benenne ihn und korrigiere die Rechnung.
\teil Bestimme die Gleichung der Geraden durch $R(1\,|\,{-1})$ und $T(3\,|\,7)$.
\end{teile}
\end{aufgabe}

\begin{aufgabe}{\textbf{Ich kann begründen, wie zwei Geraden zueinander liegen.} Entscheide ohne Zeichnen und ohne Rechnung: Schneiden sich die Geraden, sind sie parallel oder liegen sie aufeinander? Begründe.}
\begin{teile}
\teil $g$: $y = 2x - 3$ \quad und \quad $h$: $y = 2x + 4$
\teil $g$: $y = -x + 2$ \quad und \quad $h$: $y = x + 2$
\teil $g$: $y = 0{,}5x + 1$ \quad und \quad $h$: $y = 1 + 0{,}5x$
\end{teile}
\end{aufgabe}

\begin{aufgabe}{\textbf{Ich kann eine Geradengleichung im Sachzusammenhang bestimmen.} Ein Heißluftballon sinkt gleichmäßig. 4 Minuten nach Beginn der Messung ist er 820 m hoch, nach 10 Minuten nur noch 700 m. $x$ ist die Zeit in Minuten, $y$ die Höhe in m.}
\begin{teile}
\teil Bestimme die Gleichung, die die Höhe des Ballons beschreibt. (P10 2015)
\teil Was bedeuten $m$ und $n$ im Sachzusammenhang?
\teil Nach wie vielen Minuten landet der Ballon?
\end{teile}
\end{aufgabe}
